\chapter{Flux Control Design and Validation}

\section{Chapter Objective and Interface from Identification}

This chapter develops the flux-control stage based on the control-use model obtained in Chapter 3. The objective is to translate identified dynamics into practical controller design and to quantify how design choices affect bandwidth, decoupling behavior, and robustness under tested operating conditions.

\section{Flux-Control Problem Definition and Design Targets}

This section defines the inner-loop control problem, including tracking objective, decoupling requirement, disturbance sensitivity, and target dynamic behavior. These targets provide a common basis for comparing controller structures in the rest of this chapter.

\section{Baseline PI Control Design and Initial Validation}

This section presents the baseline PI design path and its expected behavior from model-based prediction. Simulation and experimental results are used together to identify where baseline design is sufficient and where performance limitations emerge.

\section{Model-Based Control with Disturbance Rejection}

This section develops the model-based control structure, including disturbance-rejection elements and design parameters that shape closed-loop behavior. The emphasis is on design rationale, not only performance display, so each control component is linked to a specific physical issue.

\section{Component-Level Interpretation and Decoupling Analysis}

This section compares controller structures at component level and interprets their effects on coupling suppression, tracking behavior, and sensitivity to model mismatch. The analysis provides the causal link required for downstream force-generation interpretation.

\section{Simulation and Experimental Validation}

This section provides paired simulation and experimental validation across representative test conditions. The results show where predicted and measured behavior are consistent and establish the validated controller outputs that will be consumed by Chapter 5.

\section{Chapter Summary and Interface to Chapter 5}

This chapter concludes with the accepted flux-control configuration and its validated behavior envelope. It then defines the interface handoff to the force-generation stage, including which outputs are treated as trusted inputs for force-quality analysis.
